% =============================================================================
% NOTA EDITORIAL — Moldura paratextual e metodologia de reconstituição
% =============================================================================
\cleardoublepage
\thispagestyle{empty}
\chapter*{Apêndice --- Nota Editorial}
\addcontentsline{toc}{chapter}{Apêndice --- Nota Editorial}
\markboth{Nota Editorial}{}

\begingroup

Esta edição foi preparada a partir da transcrição pericial do caderno original do Paciente 1.260, encontrado em anexo ao prontuário do Hospital Colônia de Barbacena sob custódia do Instituto Médico Legal em 1979.

\vspace{0.6em}

\noindent\textbf{Metodologia:}
O critério editorial adotado foi o de máxima fidelidade ao documento original, preservando sua estrutura fragmentária, suas lacunas, suas repetições e suas inconsistências internas. Os 84 fragmentos aqui reunidos seguem a ordem aproximada de escrita, reconstituída por datação relativa do suporte, da tinta e da caligrafia, conforme laudo pericial anexo (DOC-03).

\vspace{0.6em}

\noindent\textbf{Aparato paratextual:}
\begin{itemize}
  \item \textbf{[N.E.]} --- Nota do Editor: comentário editorial sobre o conteúdo, a fonte ou os limites da reconstituição.
  \item \textbf{[ilegível]} --- Trecho do original não decifrado por degradação do suporte.
  \item \textbf{[lacuna]} --- Página ou trecho ausente do caderno original.
  \item \textbf{[sic]} --- Transcrição literal de erro ou idiossincrasia do original.
\end{itemize}

\vspace{0.6em}

\noindent\textbf{Limites:}
Esta edição não reivindica completude. O diário original continha aproximadamente 847 páginas; o que chega ao leitor é uma seleção editada, organizada em cinco dimensões temáticas (MEM, DOC, LUC, CONT) mais o Epílogo. Fragmentos suprimidos por redundância, degradação irreversível ou proteção à privacidade de terceiros constam no inventário do acervo sob sigilo pericial.

\vspace{0.6em}

\noindent\textbf{Advertência final:}
O presente volume é uma reconstituição, não uma reprodução fac-similar. As escolhas editoriais --- ordenação, seleção, anotação --- refletem o estado da pesquisa na data de publicação. Pesquisadores futuros, com acesso a técnicas não destrutivas adicionais, poderão ampliar ou corrigir o que aqui se apresenta.

\vspace{1.2em}

\begin{flushright}
--- A Edição\par
OpenCode Books, Crateús\par
Julho de 2026
\end{flushright}

\endgroup
\clearpage
